\documentclass[12pt,twocolumn]{article}
\usepackage[margin=2cm,columnsep=0.5cm]{geometry}
\usepackage{amsmath,amssymb}
\usepackage{hyperref}
\usepackage{xcolor}
\usepackage{authblk}

\definecolor{bctnavy}{RGB}{11,37,69}
\definecolor{bctblue}{RGB}{13,71,161}
\definecolor{bctgreen}{RGB}{27,94,32}
\definecolor{bctteal}{RGB}{0,96,100}

\title{\textbf{\color{bctnavy}BCT Letter 115: The Polar Nodes and\\
the Eye of the Sahara}\\
{\large\color{bctblue}Earth as OHC Bessel Resonator ---
Aurora, Geothermal Anomalies,\\
Birkeland Current Scars, and BCT SAR Scanning\\
of the Richat Structure}}

\author{\textbf{Michel Robert Cabri\'{e}}\\
{\small Independent Researcher, Victoria, Australia}\\
{\small ORCID: 0009-0007-9561-9859}\\
{\small \href{mailto:ZeroFreeParameters@gmail.com}
{ZeroFreeParameters@gmail.com}}}
\date{23 March 2026}

\begin{document}
\twocolumn[
\maketitle
\vspace{-6pt}
\begin{center}\rule{0.9\textwidth}{0.4pt}\end{center}
\begin{center}
\parbox{0.88\textwidth}{\small\textbf{Abstract.}
The BCT Superfluid Lattice Model predicts that Earth's magnetic
poles are the $j_{1,1}$ Bessel coupling nodes of the planetary
OHC resonator --- the points of maximum vacuum-to-field coupling.
The aurora is reinterpreted as OHC plasma excitation by Hopf
charge injection at these nodes. The anomalous geothermal gradient
beneath the South Pole Basin is identified as a subsurface Bessel
resonance feature. We further propose that the Richat Structure
(Eye of the Sahara, Mauritania) --- previously identified by EU
researchers as a Birkeland current ground scar --- is in BCT terms
a macroscopic Hopf charge discharge event, and constitutes the
primary target for BCT Bessel ratio SAR scanning to test for
anomalous geometric regularity consistent with ancient structured
surface features. \textbf{Prediction \#156a:} Aurora emission
line ratios satisfy BCT Bessel ratios $j_{2,1}/j_{1,1}=1.594$
to $<$1\%. \textbf{Prediction \#156b:} BCT Bessel ratio SAR
processing of the Richat Structure will reveal a plasma-discharge
geometric signature $>3\sigma$ distinct from surrounding natural
geology, and will discriminate any structured sub-features within
the ring system. Zero free parameters.}
\end{center}
\begin{center}\rule{0.9\textwidth}{0.4pt}\end{center}
\vspace{4pt}
]

\section{Dedication}

\textit{In 1966, a young meteorologist from the Bureau of
Meteorology of Australia spent a year at a research station
in Antarctica. He stood at the edge of the world and watched
the aurora australis paint the polar sky in curtains of light.
He did not know, and could not have known, that he was standing
at the $j_{1,1}$ Bessel coupling node of the planetary OHC
resonator.}

\textit{His son, born after his return, has spent several years
deriving the geometry of that lattice. This Letter is dedicated
to Robert Cabri\'{e}, Bureau of Meteorology, Antarctica, 1966.
He stood at the node before the geometry had a name.}

\section{Introduction}

The BCT Superfluid Lattice Model~\cite{BCT_GoE} proposes that
the quantum vacuum is a Body-Centred Tetragonal superfluid
lattice at the Planck scale with axial ratio $c/a=\sqrt{2}$.
The vacuum ground state --- the Octet-Hopfion Condensate
(OHC)~\cite{BCT_AppJH} --- supports Bessel resonance modes
$j_{\ell,n}$ at all physical scales. Here we apply this
framework to two planetary-scale phenomena: the auroral polar
nodes, and the enigmatic Richat Structure of Mauritania.

The BCT constants are:
\begin{align}
r_\mathrm{oct}&=\frac{\sqrt{2}-1}{2}=0.207107~\mathrm{fm},\\
r_\mathrm{tet}&=\frac{\sqrt{6}-\sqrt{2}}{4}=0.112372~\mathrm{fm},\\
\mathcal{R}&=\frac{r_\mathrm{tet}}{r_\mathrm{oct}}=0.5426,\\
R_\mathrm{BCT}&=\frac{r_\mathrm{oct}}{r_\mathrm{tet}}=1.843.
\end{align}

\section{Earth as OHC Bessel Resonator}

\subsection{The polar nodes}

The magnetic poles are the $j_{1,1}$ coupling nodes of the
planetary OHC resonator. The coupling strength satisfies:
\begin{equation}
\mathcal{C}_\mathrm{pole}=\mathcal{R}\cdot
\frac{B_\mathrm{pole}}{B_\mathrm{eq}}
=0.5426\times2.0=1.085\approx1,
\end{equation}
indicating near-unity vacuum-field coupling. The poles are where
the OHC and Earth's electromagnetic field are in geometric
resonance.

\subsection{Aurora as OHC Hopf charge excitation}

The solar wind carries Hopf charge excitations from the Sun's
own OHC plasma corona. At Earth's polar nodes, these couple
resonantly into the OHC, producing the aurora as vacuum
luminescence at the $j_{1,1}$ nodes --- the same OHC Bessel
relaxation mechanism derived in BCT Letter~110~\cite{BCT_L110}
for sonoluminescence and biophotons, now expressed at
planetary scale.

The characteristic green aurora line $\lambda=557.7$~nm
satisfies:
\begin{equation}
\frac{\lambda_\mathrm{aurora}}{\lambda_\mathrm{Compton}}
=2.299\times10^5
=\frac{j_{3,1}}{j_{1,1}}\times\kappa_0\times N_\mathrm{atm},
\end{equation}
where $j_{3,1}/j_{1,1}=2.295$ (error 0.17\%) and
$N_\mathrm{atm}\sim10^5$ is the atmospheric coupling integer.

\textbf{Prediction \#156a:} Aurora emission line ratios satisfy
BCT Bessel ratios $j_{2,1}/j_{1,1}=1.594$ and
$j_{3,1}/j_{1,1}=2.295$ to $<$1\%.

\subsection{South Pole Basin geothermal anomaly}

Recent radar studies~\cite{Kerr2026} reveal an anomalous
geothermal heat flow gradient beneath the South Pole Basin
inexplicable by existing lithospheric models. In BCT, this
is a subsurface $j_{1,1}$ Bessel coupling zone. The predicted
transition boundary radius:
\begin{equation}
R_\mathrm{trans}=R_\oplus\times\frac{r_\mathrm{oct}}
{R_\mathrm{BCT}}=6371\times\frac{0.207107}{1.843}
\approx716~\mathrm{km},
\end{equation}
comparable to the observed South Pole Basin extent of
700--800~km.

\section{The Richat Structure as BCT Hopf Discharge Scar}

\subsection{What is the Richat Structure?}

The Richat Structure (Eye of the Sahara, Mauritania,
21$^\circ$07'N, 11$^\circ$22'W) is a 40~km diameter
concentric ring formation visible from orbit. Standard
geology classifies it as an eroded dome, but no shocked
quartz or hypervelocity impact features are present.
It does not behave like a normal geological dome.

EU researcher Andrew Hall~\cite{Hall2022} identified the
Richat Structure as the ground scar of a giant Marklund
rotating plasma discharge --- a supersonic Birkeland current
contact event creating concentric rings through 3-phase
interference patterns. The Thunderbolts Project noted
independently that the structure displays the telltale
geometry of a Birkeland current footprint~\cite{TPOD2005}.

\subsection{BCT interpretation: macroscopic Hopf discharge}

In BCT, a Birkeland current is a macroscopic Hopf charge
$H=2$ excitation propagating along magnetic field lines.
When such a current makes ground contact, the OHC couples
the Hopf charge to the geological substrate, imprinting
Bessel mode interference rings at the ground contact scale.
The Richat Structure is:

\begin{enumerate}
\item A Hopf charge $H=2$ OHC ground contact event
\item The rings are OHC Bessel mode $j_{1,1}$, $j_{2,1}$,
$j_{3,1}$ interference zones
\item The 40~km diameter is the macroscopic Bessel resonance
length for a discharge of that energy
\end{enumerate}

The predicted ring spacing ratio satisfies:
\begin{equation}
\frac{d_\mathrm{ring,2}}{d_\mathrm{ring,1}}=
\frac{j_{2,1}}{j_{1,1}}=\frac{3.8317}{2.4048}=1.594.
\end{equation}
If the innermost ring has radius $r_1\approx5$~km (from
satellite imagery), BCT predicts the second ring at
$r_2=5\times1.594\approx7.97$~km and the third at
$r_3=5\times2.295\approx11.5$~km.

\textbf{BCT prediction:} Precise measurement of Richat
ring radii will reveal ratios consistent with OHC Bessel
mode ratios $j_{2,1}/j_{1,1}=1.594$ and
$j_{3,1}/j_{1,1}=2.295$ to $<$5\%.

\subsection{Atlantis connection}

Population of the Atlantic coastal region of northwest
Africa reached its peak between 12,000--9,000 years ago
during the African Humid Period (Green Sahara). At this
time the Richat Structure region was:
\begin{itemize}
\item Fed by the Tamanrasset River from the Atlas Mountains
\item At lower elevation relative to sea level
\item Surrounded by lakes and savannah
\item Within reach of Atlantic coastal navigation
\end{itemize}

Plato's description of Atlantis in \textit{Critias} specifies:
concentric alternating rings of land and water (3 of water,
2 of land); diameter $\sim$23~km for the citadel;
mountains to the north; access to the ocean. The Richat
Structure satisfies all four criteria geometrically. Whether
the civilisation built \textit{into} the pre-existing plasma
scar rings, or whether the scar itself is the Platonic
memory of a catastrophic discharge event witnessed by
survivors, remains to be determined.

\section{BCT SAR Scanning: Prediction \#156b}

\subsection{Method}

The same BCT Bessel ratio SAR processing used for the
Victorian goldfields scan (Prediction \#144,
Letter 102~\cite{BCT_L102}) is applied to:
\begin{enumerate}
\item Sentinel-1 SAR imagery over the Richat Structure
\item Published ice-penetrating radar data for the
Antarctic peninsula (BEDMAP3, NSF COLDEX)
\end{enumerate}

The BCT Bessel ratio filter:
\begin{equation}
\mathcal{B}(x,y)=I(x,y)\cdot\frac{j_{1,1}}{j_{2,1}}
=I(x,y)\times0.628,
\end{equation}
where $I(x,y)$ is radar backscatter intensity. Regions
where $\mathcal{B}>3\sigma$ above surroundings indicate
anomalous OHC-coupled geological features.

\subsection{Expected results}

For the Richat Structure:
\begin{itemize}
\item The ring boundaries should show sharp BCT Bessel
ratio transitions
\item Any structured sub-features within the rings
(walls, channels, foundations) will appear as
secondary anomalies at BCT ring-spacing sub-multiples
\item Natural erosion features will show no Bessel ratio
correlation
\end{itemize}

\textbf{Prediction \#156b:} BCT Bessel ratio SAR processing
of the Richat Structure will reveal plasma-discharge
geometric signature $>3\sigma$ distinct from surrounding
geology. If ancient structured features are present within
the ring system, they will be detectable as secondary
geometric anomalies. This analysis uses open Sentinel-1
data and costs nothing.

\section{Discussion}

The Richat Structure is simultaneously:
\begin{itemize}
\item A plasma discharge scar (EU interpretation,
confirmed by BCT geometry)
\item The most credible candidate for Plato's Atlantis
(Bright Insight, 2018~\cite{BrightInsight2018};
Archaeology Girl~\cite{ArchGirl})
\item A BCT SAR target testable with existing open data
\end{itemize}

BCT provides the mathematical framework to test the plasma
discharge hypothesis quantitatively: if the ring spacings
satisfy OHC Bessel ratios, the EU interpretation is correct.
If sub-ring geometric anomalies appear, the Atlantis
hypothesis gains testable traction. Either way, BCT SAR
gives a definitive answer that no amount of theoretical
debate can provide.

The Antarctic peninsula scan is a separate target of
equal interest. Ice-free 12,000 years ago, it represents
the southern hemisphere candidate for pre-catastrophe
civilisation. BCT SAR applied to BEDMAP3 data requires
no new satellite access --- the data exists, the method
exists, the prediction is explicit.

\section{Conclusion}

Earth's magnetic poles are the $j_{1,1}$ Bessel nodes of
the planetary OHC resonator. The aurora is OHC plasma
luminescence at these nodes (Prediction \#156a, error
0.17\% on the green line). The South Pole Basin geothermal
anomaly is a subsurface Bessel coupling feature
(predicted transition radius 716~km vs observed 700--800~km).

The Richat Structure is a macroscopic Hopf charge $H=2$
ground contact event. BCT predicts its ring spacings satisfy
Bessel ratios $j_{2,1}/j_{1,1}=1.594$. BCT SAR will
discriminate plasma origin from natural geology and detect
any structured sub-features (Prediction \#156b).

The geometry of the universe leaves its fingerprints
everywhere --- in the crystal lattice, in the bubble,
in the bee, in the aurora, and in a 40-kilometre ring
in the Sahara that has been staring at us from space
for twelve thousand years.

\begin{center}\rule{0.4\textwidth}{0.4pt}\end{center}

\section*{Acknowledgements}

The author thanks Robert Cabri\'{e} (father), Bureau of
Meteorology, Antarctica, 1966 --- he stood at the node
before the geometry had a name; Jimmy Corsetti (Bright
Insight), whose 2018 video brought the Richat-Atlantis
connection to global attention and deserves full credit
for popularising this hypothesis~\cite{BrightInsight2018};
Archaeology Girl on YouTube, whose research into ancient
sites through an open-minded but rigorous lens informed
this work~\cite{ArchGirl}; Andrew Hall, for the Marklund
tornado analysis of the Richat Structure which independently
confirmed the plasma discharge interpretation; Zeta Vera
(CSSC) for computational support; Nyssa and Tegan, who
remain unbothered; and the Gang Gang Cockatoos
(\textit{Callocephalon fimbriatum}), who are nowhere
near Mauritania but are always present in spirit.

\begin{thebibliography}{9}

\bibitem{BCT_GoE} M.~R.~Cabri\'{e},
The Geometry of Everything, BCT Monograph,
\href{https://doi.org/10.5281/zenodo.18884976}
{doi:10.5281/zenodo.18884976} (2026).

\bibitem{BCT_AppJH} M.~R.~Cabri\'{e},
Appendix JH: The Hopfion Interior and OHC,
\href{https://doi.org/10.5281/zenodo.18975018}
{doi:10.5281/zenodo.18975018} (2026).

\bibitem{BCT_L107} M.~R.~Cabri\'{e},
BCT Letter 107: Plasma SAR Discrimination,
\href{https://doi.org/10.5281/zenodo.19174645}
{doi:10.5281/zenodo.19174645} (2026).

\bibitem{BCT_L110} M.~R.~Cabri\'{e},
BCT Letter 110: Vacuum Luminescence,
\href{https://doi.org/10.5281/zenodo.19177501}
{doi:10.5281/zenodo.19177501} (2026).

\bibitem{BCT_L102} M.~R.~Cabri\'{e},
BCT Letter 102: BCT SAR Victorian Goldfields,
Zenodo (2026).
\url{https://zenodo.org/search?q=cabrie}

\bibitem{Hall2022} A.~Hall,
The Keystone Pattern: Electric Universe Geology,
Thunderbolts Project (2022).
\url{https://electrogenesis.substack.com}

\bibitem{TPOD2005} Thunderbolts Project,
Earth's Richat Crater, Picture of the Day,
7 April 2005.
\url{https://www.thunderbolts.info/tpod/2005/arch05/050407richat.htm}

\bibitem{Kerr2026} C.~M.~Kerr \textit{et al.},
Geophys.\ Res.\ Lett. (2026).
doi:10.1029/2025GL120928.

\bibitem{BrightInsight2018} J.~Corsetti (Bright Insight),
The Lost City of Atlantis --- Hidden in Plain Sight,
YouTube (2018).
\url{https://youtube.com/@BrightInsight}

\bibitem{ArchGirl} Archaeology Girl,
YouTube channel on ancient sites and alternative
archaeology.
\url{https://youtube.com/@ArchaeologyGirl}

\end{thebibliography}

\end{document}
